\section{Contextualización y definición del problema}
%\begin{orientacion}
%Explique la necesidad de automatizar la limpieza de tanques, tuberías o equipos de los sectores de alimentos y bebidas. Sustente con fuentes verificables la importancia de contar con ciclos repetibles, controlados, seguros y trazables. Delimite el equipo o circuito de prueba y justifique la portabilidad de la unidad. Use citas numéricas IEEE, por ejemplo: ``los sistemas CIP reducen la intervención manual'' \cite{pi_institucional}.
%\end{orientacion}
{En el marco de una visita técnica realizada a INDUSTRIAS S.A., se identificó la oportunidad de estudiar la automatización de sistemas de limpieza para equipos empleados en los sectores de alimentos, bebidas y embotellado industrial. La empresa participa como aliado y referente del contexto industrial, evidenciando la necesidad de contar con procesos de limpieza repetibles, seguros y verificables sobre las máquinas de producción \cite{bowser2023cip} \cite{pi_institucional}. El proyecto no pretende reproducir una máquina propiedad de la empresa ni utilizar su diseño CAD; en su lugar, el equipo desarrolla una solución propia, modular y portátil, aplicable a diferentes equipos industriales del sector. La limpieza en sitio, conocida como CIP por su sigla en inglés (Clean-In-Place), permite higienizar las superficies internas de tanques, tuberías, válvulas y otros equipos de proceso sin desmontarlos completamente, haciendo circular agua o soluciones de limpieza por las mismas rutas que recorren normalmente los productos durante la operación \cite{bowser2023cip}. Frente a la limpieza manual, esta metodología facilita la repetibilidad de los ciclos, reduce la intervención del operario y favorece la trazabilidad documental exigida en procesos de manufactura de alimentos y bebidas.}

\subsection{Planteamiento del problema}
{Las máquinas empleadas en procesos industriales de alimentos, bebidas y embotellado frecuentemente carecen de una unidad de limpieza portátil que permita ejecutar de manera automática y controlada las fases de preenjuague, lavado y enjuague final directamente sobre el equipo de producción, sin depender de una instalación CIP fija dedicada a una sola línea. En ausencia de esta unidad, la limpieza de estos equipos suele realizarse de forma manual o con procedimientos no estandarizados, sin control de temperatura, sin dosificación controlada del producto de limpieza y sin registro sistemático de las variables del proceso, lo que dificulta garantizar la repetibilidad, la seguridad y la trazabilidad exigidas por el sector. Esta situación motiva la pregunta de diseño que orienta el proyecto: ¿cómo diseñar conceptualmente un sistema portátil de limpieza CIP, aplicable a distintos equipos industriales, que permita ejecutar, controlar y supervisar de forma automática un ciclo completo de limpieza, dentro de las restricciones de tiempo, presupuesto y alcance de un proyecto integrador de octavo semestre?}

\subsection{Justificación}
{Desde el punto de vista técnico, el proyecto integra de manera articulada las áreas de diseño mecatrónico, sistemas de control y comunicaciones industriales alrededor de un caso de aplicación real identificado en campo, lo que permite a los integrantes del equipo aplicar de forma conjunta conocimientos que normalmente se abordan por separado. Desde el punto de vista industrial, el proyecto responde a una necesidad concreta manifestada por INDUSTRIAS TANUZI S.A. durante la visita técnica, la cual es contar con una solución portátil, modular y transportable que pueda conectarse temporalmente a distintos equipos de producción para ejecutar ciclos de limpieza repetibles, seguros y verificables, sin exigir una instalación CIP fija por cada máquina \cite{pi_institucional}. Desde el punto de vista académico, el desarrollo de un prototipo funcional permite validar, sobre un caso industrial real, estrategias de control de temperatura, arquitecturas de comunicaciones industriales y supervisión IoT. Por otro lado, desde el punto de vista económico, una unidad portátil reduce la inversión necesaria frente a instalar sistemas CIP fijos en cada punto de uso, al poder desplazarse entre distintos equipos de la planta, y desde el punto de vista de seguridad y ambiental, el diseño contempla el manejo controlado de fluidos calientes y del producto de limpieza dosificado, así como la conducción segura del fluido hacia el drenaje al finalizar cada fase, minimizando riesgos para el operario.}

\subsection{Objetivo general}
{Diseñar conceptualmente un sistema portátil de limpieza CIP para un equipo de prueba, integrando la arquitectura mecatrónica, el lazo de control de temperatura y la arquitectura de comunicaciones industriales necesarias para ejecutar y supervisar de forma automática las fases de preenjuague, lavado y enjuague final.}

\subsection{Objetivos específicos}
\begin{enumerate}
\item{Caracterizar el proceso de limpieza CIP y formular los requisitos verificables del sistema mediante herramientas de ingeniería de requisitos (matriz de requisitos, QFD y matrices de decisión Pugh).}
\item{Proponer una arquitectura mecatrónica conceptual, que incluya diagrama de bloques, P\&ID y distribución CAD preliminar, para un Producto Mínimo Viable capaz de preparar, calentar, circular, dosificar y recuperar el fluido de limpieza.}
\item{Definir el lazo de control preliminar de temperatura del fluido de proceso y la arquitectura de comunicaciones industriales que permitan la supervisión local y remota del sistema, manteniendo el control local independiente de la capa IoT.}
\end{enumerate}

\subsection{Alcance y limitaciones}
{En este primer corte quedarán definidos el problema, los objetivos, la caracterización preliminar del proceso CIP, la matriz de requisitos, las matrices QFD y Pugh, el diseño conceptual mecatrónico, incluyendo el diagrama de bloques, P\&ID conceptual y avance del diseño CAD, el modelo dinámico preliminar del sistema térmico, la arquitectura de comunicaciones propuesta y el plan de validación. Se aplaza para los siguientes cortes la implementación física completa del sistema, la programación definitiva del PLC, la sintonización final del controlador de temperatura, la integración completa de la capa de supervisión IoT y la ejecución de las pruebas de validación con resultados experimentales. Como limitaciones se reconocen el tiempo disponible dentro del semestre, el presupuesto asignado al proyecto y la disponibilidad de hardware de laboratorio compartido entre los subgrupos del curso.}
